\documentclass[a4paper]{article}

\begin{document}

\section{Setup}
We use the same setup of the EW papers. In particular:
\begin{itemize}
\item $m_Z = 91.1876~{\rm GeV}$, $m_W = 80.398~{\rm GeV}$, $m_H =
  125~{\rm GeV}$, $m_t = 173.2~{\rm GeV}$; 
\item $G_F = 1.16639\cdot 10^{-5}~{\rm GeV}^{-2}$;
\item $\Gamma_W = 2.1054~{\rm GeV}$, $\Gamma_Z = 2.4952~{\rm GeV}$, $\Gamma_H = 4.165~\rm{MeV}$;
\item  By default, we work in the $G_{\mu}$ scheme. 
\end{itemize}
We run at the $\sqrt{s_{\rm had}} = 14~{\rm TeV}$ LHC, and use
\texttt{NNPDF31\_nnlo\_as\_0118\_luxqed} PDFs.
We include $b$-quarks and treat them as massless. Diagrams with virtual tops are
included in the RV amplitudes but not in the VV.

\noindent
We recombine leptons
and photons as follows:
\begin{enumerate}
\item we look for the $l\gamma$ pair closest in $\Delta R$, where $l$
  is either lepton or anti-lepton
\item if the closest pair has $\Delta R \le R_{l\gamma} = 0.1$, we sum the
  two momenta and call the resulting object a ``dressed lepton''. Otherwise,
  the dressed leptons are the priginal leptons.
\end{enumerate}
For now, we do not require photon-quark isolation nor lepton-lepton isolation.
We also put a cut
\begin{equation}
  300~{\rm GeV} \le m_{l\bar l} \le \sqrt{s_{\rm had}},
\end{equation}
which is computed from dressed leptons. Finally, we run with $\mu = m_{l\bar l}/2$, again
computed from dressed leptons. 

\subsection{RV amplitudes and the complex mass scheme}
At least for now, we use
the complex mass scheme for RV only for the $Z$ and $W$ mass, so we set
$\Gamma_H=\Gamma_t=0$ in the \texttt{OpenLoops} interface. For now, this
is not yet done at tree-level so we have small inconstencies in tree-level
amplitudes from our code vs tree-level amplitudes from \texttt{OpenLoops}.

\section{Fiducial volume definition and observables}
ATLAS: at least two electrons with leading/subleading $p_\perp >
30/40~\rm{GeV}$, isolated electrons with $116 \le m_{ll} \le
1500~\rm{GeV}$, $|\Delta y_{ll}| < 3.5$. No such requirement for muons.
(Muon: different isolation). Keep the lepton rapidity fixed as
$|y_l| < 2.5$. Never use pseudo-rapidity. 

Plot: $y_l$ and $p_{\perp,l}$ $y_{ee}$, $\Delta y_{ee}$ for $m_{ll}> 116$, for $m_{ll}> 300$. Also in windows $116 - 150$, $150-200$, $200 - 300$, $300-500$,
$500-1500$ for $y_{ll}$ and $\Delta y_{ll}$. CMS: up to 2 TeV, or better the
last bin is $1500-3000$.

Ratio at different energies. CMS: same for $y_{ll}$ in $120-200$, $200-1500$. 

Maybe play with isolation.

Our setup: play a bit with $\Delta R$, starting from default 0.1. Use ATLAS
as baseline for windows, but make sure we have correct under/over flow to
keep the last 1500-3000 bin. Start from 200 (CMS) [if removing Higgs bubble
  is a mess, start from 220].
Run both inclusively and with hard cuts, to
make sure we don't loose events in the tails. Run with and without $|\Delta y_{ll}|$ cut. Instead of asymmetric cuts use product cuts. But first check
how much statistics we loose $\to$ two sets of runs, one with ATLAS asym cuts
and one with product cuts. We then use $\sqrt{p_{\perp,l^+} p_{\perp,l^-}} > 35$ GeV. 

\end{document}
